% =============================================================================
% Chapter: Cross-Modal Image Registration for Infrared and Visible Imagery
% =============================================================================

\chapter{Cross-Modal Image Registration for Infrared and Visible Imagery}
\label{ch:ir-rgb-registration}

\section{Introduction}
\label{sec:reg-intro}

The fusion of infrared (IR) and visible (RGB) imagery has become a cornerstone
of modern perception systems in autonomous driving, surveillance, remote
sensing, and military applications. Because IR sensors capture thermal radiation
emitted by objects (primarily in the long-wave infrared band, 8--14\,$\mu$m)
while visible cameras record reflected light in the 0.4--0.7\,$\mu$m range, the
two modalities convey fundamentally complementary information: IR imagery
reveals thermal signatures, penetrates smoke and haze, and is invariant to
ambient illumination, whereas RGB imagery provides rich texture, color, and
fine structural detail~\cite{ma2019infrared}. However, before any downstream
fusion, detection, or recognition system can exploit this complementarity, the
two images must be brought into precise spatial correspondence---a task known as
\emph{image registration}.

This chapter provides a graduate-level treatment of the IR--RGB registration
problem. We begin with a careful analysis of why cross-modal registration is
substantially more difficult than its single-modality counterpart
(\S\ref{sec:problem-formulation}). We then survey classical
(\S\ref{sec:classical-approaches}) and deep-learning-based
(\S\ref{sec:deep-learning-approaches}) methods, describe the principal datasets
and benchmarks (\S\ref{sec:datasets}), and conclude with a discussion of
evaluation metrics (\S\ref{sec:evaluation}).


% =============================================================================
\section{The IR--RGB Registration Problem}
\label{sec:problem-formulation}

\subsection{Spectral and Radiometric Discrepancies}
\label{subsec:spectral}

The central difficulty of IR--RGB registration is the \emph{appearance gap}
that arises from the fundamentally different physical quantities measured by
each sensor. A visible camera records scene radiance that is the product of
illumination and surface reflectance, whereas a thermal infrared camera measures
emitted radiation governed by the Stefan--Boltzmann law and the surface
emissivity of objects. Consequently, the intensity relationship between the two
modalities is highly nonlinear and scene-dependent:

\begin{itemize}
  \item \textbf{Intensity reversal.} An object that is bright in IR (high
    thermal emittance) may be dark in RGB (low reflectance), and vice versa. A
    person wearing a dark jacket appears dim in visible light but radiates
    strongly in the thermal band. Conversely, a sunlit metallic surface may be
    highly reflective in the visible spectrum yet appear cool in IR due to low
    emissivity.
  \item \textbf{Illumination dependence.} Visible imagery changes drastically
    between day and night, under shadows, and under varying weather conditions,
    while IR imagery remains largely invariant to ambient illumination. This
    means that the cross-modal appearance gap is itself non-stationary.
  \item \textbf{Texture discrepancy.} Fine textures such as printed text,
    painted lane markings, or foliage patterns that produce strong gradients in
    RGB are often invisible or heavily attenuated in IR because they lack
    thermal contrast. Conversely, thermal edges caused by differential heating
    (e.g., at the boundary between asphalt and grass) may have no visible
    counterpart.
\end{itemize}

These radiometric discrepancies violate the brightness constancy assumption
that underpins most single-modality registration methods, including optical
flow~\cite{horn1981determining} and template matching~\cite{lewis1995fast}.

\subsection{Geometric Distortion and Parallax}
\label{subsec:geometric}

Beyond appearance, practical IR--RGB systems introduce geometric challenges:

\begin{itemize}
  \item \textbf{Lens distortion.} IR and visible cameras employ different
    optical elements---germanium or chalcogenide lenses for long-wave IR versus
    glass lenses for visible---with distinct focal lengths, fields of view, and
    radial/tangential distortion profiles. Even after intrinsic calibration,
    residual distortions may differ between modalities.
  \item \textbf{Parallax from displaced viewpoints.} Unless a beam-splitter
    rig is used (as in the KAIST setup~\cite{hwang2015multispectral}), the two
    sensors are physically separated, producing parallax that varies with scene
    depth. A single homography cannot rectify this parallax for scenes with
    significant depth variation; dense, depth-dependent displacement fields are
    required.
  \item \textbf{Temporal misalignment.} In vehicle-mounted or UAV-based systems,
    differences in sensor triggering, readout speed, and frame rate cause
    temporal offsets that manifest as spatial misregistration when the platform
    or objects are in motion.
\end{itemize}

Taken together, these factors mean that IR--RGB registration demands methods
that are robust to both \emph{large appearance differences} and
\emph{spatially-varying geometric deformations}---a combination that makes the
problem substantially harder than mono-modal image registration.


% =============================================================================
\section{Classical Approaches}
\label{sec:classical-approaches}

\subsection{Feature-Based Methods}
\label{subsec:feature-based}

Feature-based registration proceeds in three stages: (i)~detect keypoints in
each image, (ii)~compute local descriptors and match across modalities, and
(iii)~estimate a geometric transformation from the correspondences.

\paragraph{SIFT, SURF, and ORB.}
Scale-Invariant Feature Transform (SIFT)~\cite{lowe2004distinctive} and
Speeded-Up Robust Features (SURF)~\cite{bay2008speeded} descriptors were
designed for matching images of the same scene under viewpoint and illumination
changes within a single modality. Because they encode local gradient
orientation histograms, they implicitly assume that gradient patterns are
preserved across images. In cross-modal IR--RGB pairs, this assumption is
frequently violated: edges in IR arise from thermal boundaries whereas edges in
RGB arise from reflectance boundaries, producing inconsistent gradient fields.
ORB~\cite{rublee2011orb}, which relies on intensity comparisons (BRIEF
descriptors) and FAST corners, suffers from the same fundamental limitation.

Several adaptations have been proposed. Chen \textit{et al.}~\cite{chen2022infrared}
introduced modifications to the SIFT descriptor pipeline that emphasize
gradient magnitude rather than orientation, exploiting the observation that
strong edges (e.g., building silhouettes) tend to be co-located in both
modalities even when their contrast polarity differs. Edge-oriented
descriptors such as the histogram of oriented phase congruency
(HOPC)~\cite{ye2017robust} achieve greater modality invariance by operating on
phase congruency maps rather than raw gradients. Nevertheless, feature-based
methods remain fragile in regions that lack structural edges common to both
modalities---a frequent occurrence in natural outdoor scenes.

\paragraph{RANSAC and Transformation Estimation.}
Given a set of putative matches, Random Sample Consensus
(RANSAC)~\cite{fischler1981random} is used to robustly estimate a global
transformation---typically a homography or affine model. Because the inlier
ratio in cross-modal matching is often very low (below 10\%), RANSAC requires
many iterations and may fail entirely when too few correct matches exist.

\subsection{Intensity-Based Methods}
\label{subsec:intensity-based}

\paragraph{Mutual Information.}
Mutual information (MI) measures the statistical dependence between the
intensity distributions of two images and does not require a functional
relationship between pixel values~\cite{viola1997alignment,
maes1997multimodality}. Registration is cast as maximizing MI (or its
normalized variant, NMI) over the space of geometric transformations:
\begin{equation}
  \hat{\mathbf{T}} = \arg\max_{\mathbf{T}} \;
    \mathrm{MI}\bigl(I_{\mathrm{IR}},\; I_{\mathrm{RGB}} \circ \mathbf{T}\bigr),
  \label{eq:mi-registration}
\end{equation}
where $\mathbf{T}$ parameterizes the spatial transformation. MI-based
registration has been the workhorse of medical image registration for decades
and has been applied to IR--RGB pairs. However, MI maximization has a narrow
convergence basin, is prone to local optima, and typically operates at the
global transformation level (rigid, affine, or low-order polynomial), making it
ill-suited for recovering dense, spatially-varying deformations caused by
parallax. Furthermore, because MI is computed over intensity histograms, it
discards spatial structure and can be insensitive to local misalignments.

\paragraph{Phase Correlation.}
Phase correlation~\cite{kuglin1975phase} estimates translation (and, via
extensions, rotation and scale) by locating the peak of the normalized cross
power spectrum between two images. The method is efficient and robust to
illumination differences within a single modality. Cross-modally, however, the
spectral content of IR and RGB images can differ so dramatically that the phase
correlation peak becomes diffuse or spurious, degrading registration accuracy.

\subsection{Structural Similarity Descriptors}
\label{subsec:structural}

To bridge the appearance gap while remaining within an optimization framework,
several structural descriptors have been proposed that map each pixel to a
modality-invariant representation:

\begin{itemize}
  \item \textbf{Modality Independent Neighbourhood Descriptor (MIND)}~\cite{
    heinrich2012mind}: Encodes the local self-similarity structure around each
    voxel, producing a descriptor that is invariant to contrast differences
    between modalities. MIND and its multi-scale variant (macMIND) have been
    used as loss functions in learning-based cross-modal registration
    networks~\cite{macjnet2023}.
  \item \textbf{Self-Similarity Context (SSC)}~\cite{heinrich2013towards}:
    Extends MIND by aggregating self-similarity over a larger spatial context,
    improving discriminability at the cost of increased computation.
\end{itemize}

While these hand-crafted structural descriptors improve cross-modal robustness,
they cannot capture the full complexity of the IR--RGB appearance mapping,
motivating the transition to learned representations.


% =============================================================================
\section{Deep Learning Approaches}
\label{sec:deep-learning-approaches}

\subsection{Learning Cross-Modal Feature Descriptors}
\label{subsec:learned-descriptors}

A natural entry point for deep learning in cross-modal registration is to
replace hand-crafted descriptors with \emph{learned} modality-invariant feature
embeddings.

\paragraph{Siamese and Pseudo-Siamese Networks.}
The standard architecture for learning cross-modal descriptors is the Siamese
network, in which two branches share parameters and are trained with a
contrastive or triplet loss to map corresponding patches from different
modalities to nearby points in a shared embedding space, and non-corresponding
patches to distant points~\cite{zagoruyko2015learning}.

In the cross-modal setting, a refinement known as the \emph{pseudo-Siamese}
architecture has proven effective~\cite{cmm-net2021}. Here, the shallow layers
of each branch have \emph{unshared} parameters that learn modality-specific
low-level features (e.g., IR-specific thermal gradients vs.\ RGB-specific color
edges), while deeper layers share parameters and learn a modality-invariant
representation. This design acknowledges that the statistical properties of
the two modalities diverge most strongly at low levels of abstraction but
converge at higher semantic levels. CMM-Net~\cite{cmm-net2021} demonstrated
this principle on airborne--ground thermal--visible matching, achieving
state-of-the-art recall on cross-modal retrieval benchmarks.

Recent work has augmented Siamese backbones with Transformer
encoders~\cite{vaswani2017attention} to capture long-range dependencies in the
multi-scale feature maps produced by CNN branches, enabling attention-based
cross-modal feature comparison~\cite{deep-siamese-multilevel2024}.

\paragraph{Loss Functions for Cross-Modal Embedding.}
The choice of training loss is critical. Contrastive loss~\cite{
hadsell2006dimensionality} and triplet loss~\cite{schroff2015facenet} are
standard, but cross-modal settings benefit from additional regularization:
maximum mean discrepancy (MMD) penalties that reduce the distribution shift
between modality-specific features, and hard-negative mining strategies that
focus training on the most confusing cross-modal pairs.

\subsection{End-to-End Registration Networks}
\label{subsec:end-to-end}

Rather than detecting and matching discrete keypoints, end-to-end registration
networks directly predict a dense spatial transformation from an input image
pair.

\paragraph{Architecture.}
The dominant paradigm, pioneered by Spatial Transformer Networks
(STN)~\cite{jaderberg2015spatial} and extended by FlowNet~\cite{
dosovitskiy2015flownet} and VoxelMorph~\cite{balakrishnan2019voxelmorph},
takes a concatenated (or dual-stream encoded) image pair as input and outputs a
dense displacement field $\boldsymbol{\phi}: \Omega \to \mathbb{R}^2$ that
specifies, for each pixel $\mathbf{p}$ in the source image, its corresponding
location $\mathbf{p} + \boldsymbol{\phi}(\mathbf{p})$ in the target image. The
source image is then warped to the target coordinate frame via differentiable
bilinear interpolation, enabling end-to-end training.

In the encoder--decoder formulation, a U-Net-style backbone extracts multi-scale
features from the input pair, and the decoder progressively upsamples and
refines the predicted flow field. Skip connections between encoder and decoder
layers preserve fine-grained spatial information essential for sub-pixel
alignment accuracy.

\paragraph{Coarse-to-Fine and Cascaded Architectures.}
For large displacements---common in IR--RGB registration due to parallax---a
single-pass prediction may be insufficient. Cascaded approaches, inspired by
the Volume Tweening Network (VTN)~\cite{zhao2019unsupervised}, stack multiple
registration subnetworks: the first predicts a coarse global alignment (e.g.,
affine), the source image is warped accordingly, and subsequent subnetworks
refine the residual deformation at progressively finer scales. Li
\textit{et al.}~\cite{li2025multistage} proposed a multi-stage coarse-to-fine
network that introduces a semantic alignment module in the fine stage,
computing both inner-product and Euclidean distance between semantic feature
maps to guide geometric refinement, preventing mutual interference between
rigid and non-rigid transformation estimates.

\paragraph{Loss Functions for Cross-Modal Registration.}
When paired ground-truth deformation fields are unavailable---as is typical---
the training loss must measure alignment quality without direct supervision on
$\boldsymbol{\phi}$. Because photometric losses (L1, L2, SSIM) assume
brightness constancy, they are inappropriate for cross-modal pairs. Instead,
several alternatives are used:

\begin{itemize}
  \item \textbf{MI-based losses.} Differentiable approximations of mutual
    information, such as Mutual Information Neural Estimation
    (MINE)~\cite{belghazi2018mutual}, enable MI maximization within a
    gradient-descent training loop. DRMIME~\cite{drmime2020} demonstrated that
    MI-based losses, combined with image pyramids, yield robust cross-modal
    registration.
  \item \textbf{MIND / macMIND losses.} The sum of squared differences (SSD)
    computed on MIND descriptors rather than raw intensities provides a
    modality-invariant photometric proxy~\cite{heinrich2012mind, macjnet2023}.
  \item \textbf{Smoothness regularization.} A spatial gradient penalty on the
    predicted deformation field encourages smooth, physically plausible
    displacements:
    \begin{equation}
      \mathcal{L}_{\mathrm{smooth}} = \sum_{\mathbf{p}} \left\|
        \nabla \boldsymbol{\phi}(\mathbf{p}) \right\|^2.
    \end{equation}
  \item \textbf{Semantic and perceptual losses.} Features extracted from a
    pre-trained network (e.g., VGG~\cite{simonyan2015very}) can serve as a
    modality-bridging representation: the perceptual loss between warped-source
    and target features penalizes semantic misalignment even when pixel
    intensities are not directly comparable.
\end{itemize}

\paragraph{Key Works.}
The GCMR framework~\cite{gcmr2023} combines a generation--registration pipeline
with direct dense prediction. An Image Cross-Modality Translation Network
(ICMTN) first produces a pseudo-IR image from the visible input, converting
the cross-modal problem to a mono-modal one; a Multi-level Residual Dense
Registration Network (MRDRN) then estimates the deformation field using MI
of the misaligned images. MulFS-CAP~\cite{mulfscap2025}, published in IEEE
TPAMI, proposes a multimodal fusion-supervised cross-modality alignment
perception framework for jointly registering and fusing unregistered IR--visible
pairs, demonstrating that coupling registration with fusion supervision improves
both tasks.

\subsection{GAN-Based Cross-Modal Translation as Preprocessing}
\label{subsec:gan-translation}

An increasingly popular strategy is to \emph{translate} one modality into
the other before applying standard mono-modal registration, thereby
side-stepping the appearance gap entirely.

\paragraph{Paired Translation: Pix2Pix.}
Pix2Pix~\cite{isola2017image} trains a conditional GAN on paired
(pixel-aligned) examples using a combination of adversarial and L1 losses. When
a calibrated, beam-splitter-aligned IR--RGB dataset is available, Pix2Pix can
learn to synthesize realistic visible images from thermal inputs (or vice
versa). TIC-CGAN~\cite{tic-cgan2019} improved upon Pix2Pix with a
coarse-to-fine generator architecture for TIR image translation in traffic
scenes, producing finer texture details.

\paragraph{Unpaired Translation: CycleGAN.}
When pixel-aligned training pairs are scarce, CycleGAN~\cite{zhu2017unpaired}
leverages cycle-consistency loss to learn bidirectional IR$\leftrightarrow$RGB
translation from unpaired data. Extensions specifically targeting the IR--RGB
domain include GMA-CycleGAN~\cite{gma-cyclegan2023}, which adds gray-image
cycle consistency and mask attention to improve salient-object texture;
PCSGAN~\cite{pcsgan2020}, which combines perceptual, pixel-wise, and
adversarial losses for thermal-to-visible face translation; and
InfraGAN~\cite{infragan2022}, which uses SSIM in its architecture to enforce
structural similarity during visible-to-IR generation.

\paragraph{Translation-Then-Register Pipelines.}
Once both images reside in (approximately) the same visual domain, standard
registration methods---from SIFT+RANSAC to deep optical flow networks---become
applicable. The GCMR framework~\cite{gcmr2023} formalizes this as the
``generation--registration paradigm.'' Empirically, the approach is most
effective when the translation model produces structurally faithful outputs; if
the GAN hallucinates textures or distorts geometry, registration accuracy
degrades. Recent work by Tang \textit{et al.}~\cite{busref2025} (BusReF)
proposes a three-stage pipeline---coarse registration, fine registration, and
fusion---using one shared set of features, demonstrating that a tightly
integrated approach outperforms disjoint translation-then-register schemes.


% =============================================================================
\section{Datasets and Benchmarks}
\label{sec:datasets}

The availability of well-aligned, multi-modal image datasets is both a
prerequisite for training learning-based registration methods and a
long-standing bottleneck for the field. We review the principal benchmarks.

\subsection{KAIST Multispectral Pedestrian Dataset}
\label{subsec:kaist}

The KAIST dataset~\cite{hwang2015multispectral} comprises 95,328 color--thermal
image pairs (640$\times$512) captured at 20\,Hz from a vehicle-mounted rig. A
beam-splitter optical system produces hardware-aligned pairs by directing the
same optical path to both a visible CCD and a thermal microbolometer. The
dataset includes 103,128 pedestrian bounding-box annotations across day and
night scenes. Because the beam-splitter enforces optical co-registration, KAIST
is widely used as ground-truth for evaluating both registration and fusion
methods. However, mechanical vibrations and calibration drift introduce
residual sub-pixel misalignments, prompting Zhang \textit{et al.}~\cite{
zhang2019weakly} to release sanitized annotations that account for
misalignment. The multi-sensor platform additionally provides RGB stereo,
3D~LiDAR, and GPS/IMU data, enabling depth-aware registration research.

\subsection{FLIR ADAS Dataset}
\label{subsec:flir}

The FLIR Starter Thermal Dataset~\cite{flir2018} contains 10,228 annotated
thermal frames for automotive object detection (28,151 person, 46,692 car, and
4,457 bicycle annotations). Crucially, the IR and RGB images in the original
dataset are \emph{not aligned}---they are captured by separate cameras with
different fields of view and resolutions. Zhang \textit{et al.} subsequently
produced a manually aligned subset of 5,141 pairs (4,128 train, 1,013
validation), making the FLIR dataset useful for studying registration under
realistic misalignment conditions.

\subsection{MSRS (Multi-Spectral Road Scenarios)}
\label{subsec:msrs}

Derived from the MFNet segmentation dataset~\cite{ha2017mfnet}, the MSRS
dataset~\cite{tang2022piafusion} provides 1,444 highly aligned IR--visible
pairs (480$\times$640) after removing 125 poorly aligned pairs from the
original 1,569. The dataset is split into 715 daytime and 729 nighttime pairs,
making it particularly valuable for studying illumination-dependent
registration and fusion. The images are in color and include per-pixel semantic
segmentation labels, enabling evaluation of registration impact on downstream
segmentation tasks.

\subsection{RoadScene}
\label{subsec:roadscene}

The RoadScene dataset~\cite{xu2020u2fusion} contains 221 registered
IR--visible image pairs captured from vehicle-mounted and hand-held cameras in
both daytime and nighttime settings. While smaller than MSRS, RoadScene offers
higher resolution color images and is frequently used as a test-time benchmark
for fusion algorithms, with registration quality evaluated qualitatively through
visual inspection of fused outputs.

\subsection{TNO Image Fusion Dataset}
\label{subsec:tno}

The TNO dataset~\cite{toet2017tno} is one of the oldest and most widely cited
benchmarks for multi-spectral image fusion. It encompasses approximately 60
pairs of registered near-infrared (NIR), intensified visual, and long-wave
infrared (LWIR) images spanning military and civilian scenarios (urban, rural,
night). The images are grayscale, relatively low resolution, and carefully
registered. While not large enough for training deep networks, TNO remains a
standard qualitative evaluation benchmark.

\subsection{Other Notable Datasets}
\label{subsec:other-datasets}

\textbf{LLVIP}~\cite{jia2021llvip} provides 16,836 aligned visible--infrared
image pairs for pedestrian detection in low-light conditions.
\textbf{CVC-14}~\cite{gonzalez2016pedestrian} offers color--thermal pairs for
pedestrian detection but suffers from persistent alignment issues due to
hardware limitations. \textbf{DroneVehicle}~\cite{sun2022drone} extends the
multi-modal paradigm to UAV-based vehicle detection with IR--RGB pairs from
aerial viewpoints.


% =============================================================================
\section{Evaluation Metrics}
\label{sec:evaluation}

Evaluating cross-modal registration quality is non-trivial because the
``correct'' alignment is often known only approximately. We categorize metrics
by the type of ground truth they require.

\subsection{Endpoint Error (EPE)}
\label{subsec:epe}

When ground-truth dense deformation fields $\boldsymbol{\phi}^*$ are available
(e.g., from synthetic warping experiments), the average endpoint error quantifies
per-pixel displacement accuracy:
\begin{equation}
  \mathrm{EPE} = \frac{1}{|\Omega|} \sum_{\mathbf{p} \in \Omega}
    \left\| \boldsymbol{\phi}(\mathbf{p})
    - \boldsymbol{\phi}^*(\mathbf{p}) \right\|_2.
  \label{eq:epe}
\end{equation}
EPE is the gold-standard metric for dense registration but requires
ground-truth flow, which is rarely available for real cross-modal data.

\subsection{Landmark-Based Metrics}
\label{subsec:landmark}

A practical alternative is to manually annotate corresponding landmarks (e.g.,
building corners, body joints) in both images and measure the Target
Registration Error (TRE):
\begin{equation}
  \mathrm{TRE} = \frac{1}{N} \sum_{i=1}^{N}
    \left\| \mathbf{p}_i^{\mathrm{warped}} - \mathbf{p}_i^{\mathrm{target}}
    \right\|_2,
\end{equation}
where $\mathbf{p}_i^{\mathrm{warped}}$ is the landmark position in the warped
source and $\mathbf{p}_i^{\mathrm{target}}$ is the corresponding position in
the target. TRE is widely used in medical registration and is increasingly
adopted for IR--RGB evaluation. Its limitation is that it measures alignment
only at sparse locations and is subject to annotation noise.

\subsection{Image Similarity After Warping}
\label{subsec:similarity}

When ground-truth deformation is unavailable, registration quality is assessed
indirectly by measuring the similarity between the warped source and the
target:

\begin{itemize}
  \item \textbf{Structural Similarity Index (SSIM)}~\cite{wang2004image}:
    Evaluates luminance, contrast, and structural agreement within local
    windows. SSIM is partially robust to cross-modal intensity differences but
    can be misleading when the modalities have fundamentally different
    appearance.
  \item \textbf{Mutual Information (MI)}: Higher MI between the warped source
    and target indicates better statistical alignment. MI is modality-agnostic
    and thus more principled than SSIM for cross-modal evaluation, though it
    measures global statistical dependence rather than spatial alignment per se.
  \item \textbf{Normalized Cross-Correlation (NCC)}: Measures linear
    correlation within local windows. NCC is effective when the intensity
    relationship is approximately affine (e.g., after GAN-based domain
    translation) but less so for raw IR--RGB pairs.
  \item \textbf{Mean Squared Error (MSE)}: Rarely appropriate for raw
    cross-modal evaluation due to the brightness constancy violation, but
    useful after domain translation when both images share the same intensity
    distribution.
\end{itemize}

\subsection{Mask-Based and Semantic Metrics}
\label{subsec:mask-based}

When per-pixel semantic labels are available for both modalities (as in MSRS),
registration can be evaluated by warping the source segmentation mask and
comparing it to the target mask:

\begin{itemize}
  \item \textbf{Intersection over Union (IoU)}: The ratio of the intersection
    to the union of the warped-source and target segmentation masks for each
    semantic class. Mean IoU (mIoU) across classes provides a holistic
    registration quality measure that is invariant to appearance differences.
  \item \textbf{Dice Coefficient}: Equivalent to the F1 score of the overlap
    between masks, closely related to IoU and widely used in medical
    registration evaluation.
\end{itemize}

These metrics have the advantage of being entirely modality-agnostic: they
measure whether \emph{semantic content} is correctly aligned regardless of
how it appears in each sensor. Their limitation is the requirement for dense
semantic annotations.

\subsection{Downstream Task Performance}
\label{subsec:downstream}

Ultimately, registration is a means to an end. An increasingly popular
evaluation strategy is to measure the impact of registration quality on a
downstream task such as pedestrian detection (miss rate on KAIST), semantic
segmentation (mIoU on MSRS), or image fusion quality (assessed via
$Q_{AB/F}$, visual information fidelity, or human perceptual
studies)~\cite{tang2022piafusion}. This task-oriented evaluation is arguably
the most meaningful but confounds registration quality with the performance
of the downstream model.


% =============================================================================
\section{Discussion and Open Challenges}
\label{sec:discussion}

Despite rapid progress, several open challenges remain in IR--RGB registration:

\begin{enumerate}
  \item \textbf{Dense deformation under parallax.} Most current methods assume
    either a global transformation or moderate local deformation. True
    depth-dependent parallax correction requires either explicit depth
    estimation or very flexible deformation models, and remains underexplored
    in the IR--RGB context.
  \item \textbf{Temporal coherence.} For video-rate IR--RGB streams, exploiting
    temporal consistency across frames can improve registration robustness, but
    few methods incorporate recurrent or spatio-temporal architectures.
  \item \textbf{Joint registration--fusion optimization.} Methods such as
    MURF~\cite{xu2023murf}, MulFS-CAP~\cite{mulfscap2025}, and
    BusReF~\cite{busref2025} demonstrate that coupling registration with fusion
    in a joint optimization yields benefits for both tasks, but the optimal
    architecture and training strategy for such joint learning remain open
    questions.
  \item \textbf{Unsupervised and self-supervised training.} The scarcity of
    ground-truth deformation fields for real IR--RGB data makes unsupervised
    approaches essential. Advances in differentiable MI estimation, structural
    descriptor losses, and self-supervised objectives (e.g., bidirectional
    self-registration~\cite{bi-directional-selfreg2025}) are promising but
    have not yet matched the accuracy of supervised methods on benchmarks with
    synthetic ground truth.
  \item \textbf{Evaluation standardization.} The field lacks a universally
    accepted benchmark protocol. Different papers evaluate on different
    subsets of different datasets with different metrics, making fair comparison
    difficult. A community effort toward standardized evaluation, akin to the
    Middlebury or KITTI benchmarks for optical flow, would significantly
    advance the field.
\end{enumerate}


% =============================================================================
% Bibliography keys used in this chapter (for the .bib file):
%
% @article{ma2019infrared,
%   title={Infrared and visible image fusion methods and applications: A survey},
%   author={Ma, Jiayi and Ma, Yong and Li, Chang},
%   journal={Information Fusion},
%   volume={45},
%   pages={153--178},
%   year={2019}
% }
%
% @article{horn1981determining,
%   title={Determining optical flow},
%   author={Horn, Berthold KP and Schunck, Brian G},
%   journal={Artificial Intelligence},
%   volume={17},
%   pages={185--203},
%   year={1981}
% }
%
% @article{lewis1995fast,
%   title={Fast normalized cross-correlation},
%   author={Lewis, JP},
%   journal={Vision Interface},
%   year={1995}
% }
%
% @article{lowe2004distinctive,
%   title={Distinctive image features from scale-invariant keypoints},
%   author={Lowe, David G},
%   journal={International Journal of Computer Vision},
%   volume={60},
%   number={2},
%   pages={91--110},
%   year={2004}
% }
%
% @article{bay2008speeded,
%   title={Speeded-up robust features ({SURF})},
%   author={Bay, Herbert and Ess, Andreas and Tuytelaars, Tinne and Van Gool, Luc},
%   journal={Computer Vision and Image Understanding},
%   volume={110},
%   number={3},
%   pages={346--359},
%   year={2008}
% }
%
% @inproceedings{rublee2011orb,
%   title={{ORB}: An efficient alternative to {SIFT} or {SURF}},
%   author={Rublee, Ethan and Rabaud, Vincent and Konolige, Kurt and Bradski, Gary},
%   booktitle={ICCV},
%   year={2011}
% }
%
% @article{chen2022infrared,
%   title={Infrared and visible image registration based on content},
%   author={Chen, Yanzhi and others},
%   journal={Pattern Recognition},
%   year={2022}
% }
%
% @article{ye2017robust,
%   title={Robust registration of multimodal remote sensing images based on structural similarity},
%   author={Ye, Yuanxin and Shan, Jie and Bruzzone, Lorenzo and Shen, Li},
%   journal={IEEE Transactions on Geoscience and Remote Sensing},
%   volume={55},
%   number={5},
%   pages={2941--2958},
%   year={2017}
% }
%
% @article{fischler1981random,
%   title={Random sample consensus: A paradigm for model fitting},
%   author={Fischler, Martin A and Bolles, Robert C},
%   journal={Communications of the ACM},
%   volume={24},
%   number={6},
%   pages={381--395},
%   year={1981}
% }
%
% @inproceedings{viola1997alignment,
%   title={Alignment by maximization of mutual information},
%   author={Viola, Paul and Wells III, William M},
%   booktitle={International Journal of Computer Vision},
%   volume={24},
%   number={2},
%   pages={137--154},
%   year={1997}
% }
%
% @article{maes1997multimodality,
%   title={Multimodality image registration by maximization of mutual information},
%   author={Maes, Frederik and Collignon, Andre and Vandermeulen, Dirk and Marchal, Guy and Suetens, Paul},
%   journal={IEEE Transactions on Medical Imaging},
%   volume={16},
%   number={2},
%   pages={187--198},
%   year={1997}
% }
%
% @inproceedings{kuglin1975phase,
%   title={The phase correlation image alignment method},
%   author={Kuglin, CD and Hines, DC},
%   booktitle={IEEE International Conference on Cybernetics and Society},
%   year={1975}
% }
%
% @inproceedings{heinrich2012mind,
%   title={{MIND}: Modality independent neighbourhood descriptor for multi-modal deformable registration},
%   author={Heinrich, Mattias P and Jenkinson, Mark and Bhushan, Manav and Matin, Tahreema and Gleeson, Fergus V and Brady, Sir Michael and Schnabel, Julia A},
%   booktitle={Medical Image Analysis},
%   volume={16},
%   number={7},
%   pages={1423--1435},
%   year={2012}
% }
%
% @inproceedings{heinrich2013towards,
%   title={Towards realtime multimodal fusion for image-guided interventions using self-similarities},
%   author={Heinrich, Mattias P and others},
%   booktitle={MICCAI},
%   year={2013}
% }
%
% @article{macjnet2023,
%   title={macJNet: weakly-supervised multimodal image deformable registration using joint learning},
%   author={Wang, Yuxi and others},
%   journal={BioMedical Engineering OnLine},
%   year={2023}
% }
%
% @inproceedings{zagoruyko2015learning,
%   title={Learning to compare image patches via convolutional neural networks},
%   author={Zagoruyko, Sergey and Komodakis, Nikos},
%   booktitle={CVPR},
%   year={2015}
% }
%
% @article{cmm-net2021,
%   title={Cross-modality image matching network with modality-invariant feature representation},
%   author={Jiang, Xin and others},
%   journal={IEEE Transactions on Geoscience and Remote Sensing},
%   year={2021}
% }
%
% @inproceedings{vaswani2017attention,
%   title={Attention is all you need},
%   author={Vaswani, Ashish and others},
%   booktitle={NeurIPS},
%   year={2017}
% }
%
% @inproceedings{deep-siamese-multilevel2024,
%   title={Deep Siamese multi-level feature network for {VIS} and {NIR} image matching},
%   author={Zhang, Wei and others},
%   booktitle={ACM International Conference on Signal Processing},
%   year={2024}
% }
%
% @inproceedings{hadsell2006dimensionality,
%   title={Dimensionality reduction by learning an invariant mapping},
%   author={Hadsell, Raia and Chopra, Sumit and LeCun, Yann},
%   booktitle={CVPR},
%   year={2006}
% }
%
% @inproceedings{schroff2015facenet,
%   title={{FaceNet}: A unified embedding for face recognition and clustering},
%   author={Schroff, Florian and Kalenichenko, Dmitry and Philbin, James},
%   booktitle={CVPR},
%   year={2015}
% }
%
% @inproceedings{jaderberg2015spatial,
%   title={Spatial transformer networks},
%   author={Jaderberg, Max and Simonyan, Karen and Zisserman, Andrew and Kavukcuoglu, Koray},
%   booktitle={NeurIPS},
%   year={2015}
% }
%
% @inproceedings{dosovitskiy2015flownet,
%   title={{FlowNet}: Learning optical flow with convolutional networks},
%   author={Dosovitskiy, Alexei and others},
%   booktitle={ICCV},
%   year={2015}
% }
%
% @article{balakrishnan2019voxelmorph,
%   title={{VoxelMorph}: A learning framework for deformable medical image registration},
%   author={Balakrishnan, Guha and Zhao, Amy and Sabuncu, Mert R and Guttag, John and Dalca, Adrian V},
%   journal={IEEE Transactions on Medical Imaging},
%   volume={38},
%   number={8},
%   pages={1788--1800},
%   year={2019}
% }
%
% @inproceedings{zhao2019unsupervised,
%   title={Unsupervised 3{D} end-to-end medical image registration with volume tweening network},
%   author={Zhao, Shengyu and Lau, Tingfung and Luo, Ji and Chang, Eric I-Chao and Xu, Yan},
%   booktitle={IEEE Journal of Biomedical and Health Informatics},
%   year={2019}
% }
%
% @article{li2025multistage,
%   title={A multi stage coarse-to-fine network with semantic similarity for infrared and visible image registration},
%   author={Li, Jun and others},
%   journal={Optics and Lasers in Engineering},
%   year={2025}
% }
%
% @article{simonyan2015very,
%   title={Very deep convolutional networks for large-scale image recognition},
%   author={Simonyan, Karen and Zisserman, Andrew},
%   journal={ICLR},
%   year={2015}
% }
%
% @article{belghazi2018mutual,
%   title={Mutual information neural estimation},
%   author={Belghazi, Mohamed Ishmael and others},
%   booktitle={ICML},
%   year={2018}
% }
%
% @article{drmime2020,
%   title={{DRMIME}: Differentiable mutual information minimization estimator for multimodal registration},
%   author={Pielawski, Nicolas and others},
%   journal={arXiv preprint arXiv:1912.05375},
%   year={2020}
% }
%
% @article{gcmr2023,
%   title={General cross-modality registration framework for visible and infrared {UAV} target image registration},
%   author={Chen, Zhiyuan and others},
%   journal={Scientific Reports},
%   year={2023}
% }
%
% @article{mulfscap2025,
%   title={Multimodal fusion-supervised cross-modality alignment perception},
%   author={Zhao, Fan and others},
%   journal={IEEE Transactions on Pattern Analysis and Machine Intelligence},
%   year={2025}
% }
%
% @inproceedings{isola2017image,
%   title={Image-to-image translation with conditional adversarial networks},
%   author={Isola, Phillip and Zhu, Jun-Yan and Zhou, Tinghui and Efros, Alexei A},
%   booktitle={CVPR},
%   year={2017}
% }
%
% @article{tic-cgan2019,
%   title={{TIC-CGAN}: Thermal image to color image translation using conditional {GAN}},
%   author={Kniaz, Vladimir V and others},
%   journal={arXiv preprint},
%   year={2019}
% }
%
% @inproceedings{zhu2017unpaired,
%   title={Unpaired image-to-image translation using cycle-consistent adversarial networks},
%   author={Zhu, Jun-Yan and Park, Taesung and Isola, Phillip and Efros, Alexei A},
%   booktitle={ICCV},
%   year={2017}
% }
%
% @article{gma-cyclegan2023,
%   title={An unpaired thermal infrared image translation method using {GMA-CycleGAN}},
%   author={Wang, Shihao and others},
%   journal={Remote Sensing},
%   year={2023}
% }
%
% @article{pcsgan2020,
%   title={{PCSGAN}: Perceptual cyclic-synthesized generative adversarial networks for thermal and {NIR} to visible image transformation},
%   author={Kancharagunta, Kishore Babu and Shyam Lal, Sonia},
%   journal={Neurocomputing},
%   year={2020}
% }
%
% @article{infragan2022,
%   title={{InfraGAN}: A {GAN} architecture to transfer visible images to infrared domain},
%   author={Oezcan, Ahmet Haydar and others},
%   journal={Pattern Recognition Letters},
%   year={2022}
% }
%
% @article{busref2025,
%   title={{BusReF}: Infrared-visible images registration and fusion},
%   author={Luo, Yuqing and others},
%   journal={ACM Transactions on Multimedia Computing},
%   year={2025}
% }
%
% @inproceedings{hwang2015multispectral,
%   title={Multispectral pedestrian detection: Benchmark dataset and baseline},
%   author={Hwang, Soonmin and Park, Jaesik and Kim, Namil and Choi, Yukyung and Kweon, In So},
%   booktitle={CVPR},
%   year={2015}
% }
%
% @article{zhang2019weakly,
%   title={Weakly aligned cross-modal learning for multispectral pedestrian detection},
%   author={Zhang, Lu and others},
%   journal={ICCV},
%   year={2019}
% }
%
% @misc{flir2018,
%   title={{FLIR} thermal dataset for algorithm training},
%   author={{FLIR Systems}},
%   year={2018},
%   howpublished={\url{https://www.flir.com/oem/adas/adas-dataset-form/}}
% }
%
% @article{ha2017mfnet,
%   title={{MFNet}: Towards real-time semantic segmentation for autonomous vehicles with multi-spectral scenes},
%   author={Ha, Qishen and Watanabe, Kohei and Karasawa, Takumi and Ushiku, Yoshitaka and Harada, Tatsuya},
%   journal={IROS},
%   year={2017}
% }
%
% @article{tang2022piafusion,
%   title={{PIAFusion}: A progressive infrared and visible image fusion network based on illumination aware},
%   author={Tang, Linfeng and Yuan, Jiteng and Zhang, Hao and Jiang, Xingyu and Ma, Jiayi},
%   journal={Information Fusion},
%   year={2022}
% }
%
% @article{xu2020u2fusion,
%   title={{U2Fusion}: A unified unsupervised image fusion network},
%   author={Xu, Han and Ma, Jiayi and Jiang, Junjun and Guo, Xiaojie and Ling, Haibin},
%   journal={IEEE Transactions on Pattern Analysis and Machine Intelligence},
%   year={2020}
% }
%
% @misc{toet2017tno,
%   title={{TNO} image fusion dataset},
%   author={Toet, Alexander},
%   year={2017},
%   howpublished={\url{https://figshare.com/articles/dataset/TNO_Image_Fusion_Dataset/1008029}}
% }
%
% @article{jia2021llvip,
%   title={{LLVIP}: A visible-infrared paired dataset for low-light vision},
%   author={Jia, Xinyu and Zhu, Ce and Li, Mingjian and Tang, Wei and Zhou, Wenli},
%   journal={ICCV Workshops},
%   year={2021}
% }
%
% @article{gonzalez2016pedestrian,
%   title={Pedestrian detection at day/night time with visible and {FIR} cameras},
%   author={Gonzalez, Alejandro and others},
%   journal={Sensors},
%   year={2016}
% }
%
% @article{sun2022drone,
%   title={{Drone-based RGB-infrared} cross-modality vehicle detection via uncertainty-aware learning},
%   author={Sun, Yiming and others},
%   journal={IEEE Transactions on Circuits and Systems for Video Technology},
%   year={2022}
% }
%
% @article{wang2004image,
%   title={Image quality assessment: From error visibility to structural similarity},
%   author={Wang, Zhou and Bovik, Alan C and Sheikh, Hamid R and Simoncelli, Eero P},
%   journal={IEEE Transactions on Image Processing},
%   volume={13},
%   number={4},
%   pages={600--612},
%   year={2004}
% }
%
% @article{xu2023murf,
%   title={{MURF}: Mutually reinforcing multi-modal image registration and fusion},
%   author={Xu, Han and others},
%   journal={IEEE Transactions on Pattern Analysis and Machine Intelligence},
%   year={2023}
% }
%
% @article{bi-directional-selfreg2025,
%   title={Bi-directional self-registration for misaligned infrared-visible image fusion},
%   author={Li, Gang and others},
%   journal={arXiv preprint arXiv:2505.06920},
%   year={2025}
% }
% =============================================================================
